\documentclass{article}
\usepackage{mdposter}
\posttitle{Predicting clone composition from relatives}
\postfield{Biology, single-cell lineage tracing}

\begin{document}
\begin{mdposter}

\posterbox{name=c1, column=1, between=top and bottom}{%
\begin{mdcol}
\mdsec{The question}
Sequencing a cell destroys it, so nobody can measure the same cell twice. In a lineage-linked assay every cell
carries an inherited barcode, so the descendants of one ancestor can be recognised as relatives afterwards. We ask when
measuring some of a clone's cells helps predict the state of its other cells, and what a rule that fails is missing:
fitting within the study, reuse of a rule from another study, or simply more cells captured.

\vspace{0.16in}
The practical question underneath is whether measuring a relative is worth what it costs. Every cell spent on
a relative is a cell not spent elsewhere, so a rule has to beat the simplest alternative, which is to ignore
the relative and predict the same number for every clone.
\mdgap
\mdsec{How we are approaching it}
\begin{itemize}
\item We score a fixed expression signature in four wells and move a two-bin predictor between the two control wells without retuning it.
\item In two barcoding data sets, of 657 and 1,801 clones, we predict a clone's composition in one group from the other, comparing a rule fitted here, a rule frozen from the other study, and simply copying the relative's value.
\item We cap the relatives measured per clone at 1, 3, 5, 10 or all of them, with exact formulas for the expected loss.
\item We refit on 10 to 400 training clones and keep between 5\% and all of the existing annotation.
\item We compare every rule against a constant fitted on the same clones, so a gain means the relative carried information the constant did not.
\end{itemize}
\end{mdcol}
}

\posterbox{name=c2, column=2, between=top and bottom}{%
\begin{mdcol}
\mdsec{The signature check}
\includegraphics[width=\linewidth]{lineage-wells.pdf}
\figcap{Figure 1.}{Where the project started. In four wells we asked how much more often a cell is in the state of interest when a relative is, against cells paired at random. All four gaps fall under the five-point gate we set before looking.}
\mdgap
\mdsec{What we have so far}
\begin{progress}
With one relative measured per clone, a rule fitted within the study beats a constant by 11.9\% and 11.2\% in the
two directions. Copying the relative's value across has a negative mean gain at every cap.
\end{progress}
\begin{progress}
All 24 frozen configurations carried over from the other study lose to a constant fitted on the target, although
each of them beats the constant from their own source.
\end{progress}
\begin{progress}
Fitting the rule needs clones. With 10 training clones per fold the mean gain is negative, and it turns positive
at about 19.
\end{progress}
\end{mdcol}
}

\posterbox{name=c3, column=3, between=top and bottom}{%
\begin{mdcol}
\mdsec{What calibration buys}
\includegraphics[width=\linewidth]{lineage-gain.pdf}
\figcap{Figure 2.}{How much each rule beats predicting a constant, in squared error, as more relatives are measured per clone. A rule fitted within the study helps from one relative onward. Copying the relative's value never helps on average, and rules frozen from the other study span the shaded band.}
\mdgap
\mdsec{What would settle it}
\begin{itemize}
\item Independent cultures with functional labels, and capture recorded as the experiment runs, rather than a further search for a signature.
\item Repeated measurements of the same state, with a model of capture, to fix the noise floor.
\end{itemize}
\mdgap
\mdsec{Open questions}
\begin{itemize}
\item Because the original capture is unknown and the cultures are not independent, none of this supports a claim about population risk or about why a treatment works.
\item We cannot say how much of the remaining error is biological rather than technical.
\item Nothing here shows an assay on one cell, or where the sequencing effort is best spent.
\end{itemize}
\end{mdcol}
}

\end{mdposter}
\end{document}
